%%%%%%%%%%%%%%%%%%%%%%%%%%%%%%%%%%%%%%%%%%%%%%%%%%%%%%%%%%%%
% 14.14 STATISTICAL TABLES
% The Composition of Advanced Mathematics
%%%%%%%%%%%%%%%%%%%%%%%%%%%%%%%%%%%%%%%%%%%%%%%%%%%%%%%%%%%%

\section{Statistical Tables}
\label{sec:statistical-tables}

\prerequisite{
Descriptive statistics, probability distributions,
sampling distributions, estimation, confidence intervals,
hypothesis testing, correlation, regression, analysis of variance,
and statistical notation.
}

\learningobjectives{
By the end of this reference section, the reader should be able to:
\begin{itemize}
    \item compute standard descriptive statistics;
    \item distinguish population and sample quantities;
    \item calculate sample variance and standard deviation;
    \item compute weighted means;
    \item standardize observations using \(z\)-scores;
    \item calculate covariance and correlation;
    \item use standard-error formulas;
    \item recognize common sampling distributions;
    \item construct standard confidence intervals;
    \item formulate common test statistics;
    \item distinguish one- and two-sided tests;
    \item interpret \(p\)-values and significance levels;
    \item use \(z\), \(t\), chi-square, and \(F\) critical values;
    \item calculate simple linear regression estimates;
    \item calculate \(R^2\);
    \item recognize ANOVA decomposition formulas;
    \item compute common proportions and contingency-table statistics;
    \item understand effect-size formulas;
    \item recognize multiple-testing corrections;
    \item avoid common statistical interpretation errors.
\end{itemize}
}

%%%%%%%%%%%%%%%%%%%%%%%%%%%%%%%%%%%%%%%%%%%%%%%%%%%%%%%%%%%%
\subsection{Population and Sample Notation}
\label{subsec:population-sample-notation}
%%%%%%%%%%%%%%%%%%%%%%%%%%%%%%%%%%%%%%%%%%%%%%%%%%%%%%%%%%%%

Common population notation:

\[
\boxed{
N
=
\text{population size},
}
\]

\[
\boxed{
\mu
=
\text{population mean},
}
\]

\[
\boxed{
\sigma^2
=
\text{population variance},
}
\]

\[
\boxed{
\sigma
=
\text{population standard deviation},
}
\]

\[
\boxed{
p
=
\text{population proportion}.
}
\]

Common sample notation:

\[
\boxed{
n
=
\text{sample size},
}
\]

\[
\boxed{
\bar{x}
=
\text{sample mean},
}
\]

\[
\boxed{
s^2
=
\text{sample variance},
}
\]

\[
\boxed{
s
=
\text{sample standard deviation},
}
\]

\[
\boxed{
\hat p
=
\text{sample proportion}.
}
\]

%%%%%%%%%%%%%%%%%%%%%%%%%%%%%%%%%%%%%%%%%%%%%%%%%%%%%%%%%%%%
\subsection{Sample Mean}
\label{subsec:sample-mean-reference}
%%%%%%%%%%%%%%%%%%%%%%%%%%%%%%%%%%%%%%%%%%%%%%%%%%%%%%%%%%%%

For observations

\[
x_1,\ldots,x_n,
\]

the sample mean is

\[
\boxed{
\bar{x}
=
\frac1n
\sum_{i=1}^{n}x_i.
}
\]

%%%%%%%%%%%%%%%%%%%%%%%%%%%%%%%%%%%%%%%%%%%%%%%%%%%%%%%%%%%%
\subsection{Population Mean}
\label{subsec:population-mean-reference}
%%%%%%%%%%%%%%%%%%%%%%%%%%%%%%%%%%%%%%%%%%%%%%%%%%%%%%%%%%%%

For finite population values

\[
x_1,\ldots,x_N,
\]

\[
\boxed{
\mu
=
\frac1N
\sum_{i=1}^{N}x_i.
}
\]

%%%%%%%%%%%%%%%%%%%%%%%%%%%%%%%%%%%%%%%%%%%%%%%%%%%%%%%%%%%%
\subsection{Weighted Mean}
\label{subsec:weighted-mean-reference}
%%%%%%%%%%%%%%%%%%%%%%%%%%%%%%%%%%%%%%%%%%%%%%%%%%%%%%%%%%%%

For weights \(w_i\),

\[
\boxed{
\bar{x}_w
=
\frac{
\sum_iw_ix_i
}{
\sum_iw_i
}.
}
\]

If

\[
\sum_iw_i=1,
\]

then

\[
\boxed{
\bar{x}_w
=
\sum_iw_ix_i.
}
\]

%%%%%%%%%%%%%%%%%%%%%%%%%%%%%%%%%%%%%%%%%%%%%%%%%%%%%%%%%%%%
\subsection{Sample Variance}
\label{subsec:sample-variance-reference}
%%%%%%%%%%%%%%%%%%%%%%%%%%%%%%%%%%%%%%%%%%%%%%%%%%%%%%%%%%%%

\[
\boxed{
s^2
=
\frac{
1
}{
n-1
}
\sum_{i=1}^{n}
(x_i-\bar{x})^2.
}
\]

The denominator

\[
n-1
\]

is the sample degrees of freedom for the usual unbiased variance estimator
under independent sampling.

%%%%%%%%%%%%%%%%%%%%%%%%%%%%%%%%%%%%%%%%%%%%%%%%%%%%%%%%%%%%
\subsection{Sample Standard Deviation}
\label{subsec:sample-standard-deviation-reference}
%%%%%%%%%%%%%%%%%%%%%%%%%%%%%%%%%%%%%%%%%%%%%%%%%%%%%%%%%%%%

\[
\boxed{
s
=
\sqrt{s^2}.
}
\]

%%%%%%%%%%%%%%%%%%%%%%%%%%%%%%%%%%%%%%%%%%%%%%%%%%%%%%%%%%%%
\subsection{Population Variance}
\label{subsec:population-variance-reference}
%%%%%%%%%%%%%%%%%%%%%%%%%%%%%%%%%%%%%%%%%%%%%%%%%%%%%%%%%%%%

\[
\boxed{
\sigma^2
=
\frac1N
\sum_{i=1}^{N}
(x_i-\mu)^2.
}
\]

%%%%%%%%%%%%%%%%%%%%%%%%%%%%%%%%%%%%%%%%%%%%%%%%%%%%%%%%%%%%
\subsection{Computational Variance Identity}
\label{subsec:computational-variance-identity}
%%%%%%%%%%%%%%%%%%%%%%%%%%%%%%%%%%%%%%%%%%%%%%%%%%%%%%%%%%%%

For a population/random variable,

\[
\boxed{
\operatorname{Var}(X)
=
\mathbb{E}[X^2]
-
(\mathbb{E}[X])^2.
}
\]

For sample calculations,

\[
\boxed{
\sum_{i=1}^{n}
(x_i-\bar{x})^2
=
\sum_{i=1}^{n}x_i^2
-
n\bar{x}^2.
}
\]

%%%%%%%%%%%%%%%%%%%%%%%%%%%%%%%%%%%%%%%%%%%%%%%%%%%%%%%%%%%%
\subsection{Range}
\label{subsec:range-statistics-reference}
%%%%%%%%%%%%%%%%%%%%%%%%%%%%%%%%%%%%%%%%%%%%%%%%%%%%%%%%%%%%

\[
\boxed{
\text{Range}
=
x_{\max}
-
x_{\min}.
}
\]

%%%%%%%%%%%%%%%%%%%%%%%%%%%%%%%%%%%%%%%%%%%%%%%%%%%%%%%%%%%%
\subsection{Interquartile Range}
\label{subsec:iqr-reference}
%%%%%%%%%%%%%%%%%%%%%%%%%%%%%%%%%%%%%%%%%%%%%%%%%%%%%%%%%%%%

\[
\boxed{
\operatorname{IQR}
=
Q_3-Q_1.
}
\]

A common exploratory outlier rule flags observations below

\[
\boxed{
Q_1-1.5\,\operatorname{IQR}
}
\]

or above

\[
\boxed{
Q_3+1.5\,\operatorname{IQR}.
}
\]

This is a descriptive convention, not a universal definition of an outlier.

%%%%%%%%%%%%%%%%%%%%%%%%%%%%%%%%%%%%%%%%%%%%%%%%%%%%%%%%%%%%
\subsection{\(z\)-Score}
\label{subsec:z-score-reference}
%%%%%%%%%%%%%%%%%%%%%%%%%%%%%%%%%%%%%%%%%%%%%%%%%%%%%%%%%%%%

For population mean \(\mu\) and standard deviation \(\sigma\),

\[
\boxed{
z
=
\frac{
x-\mu
}{
\sigma
}.
}
\]

For sample-based standardization,

\[
\boxed{
z_i
=
\frac{
x_i-\bar{x}
}{
s
}.
}
\]

%%%%%%%%%%%%%%%%%%%%%%%%%%%%%%%%%%%%%%%%%%%%%%%%%%%%%%%%%%%%
\subsection{Sample Proportion}
\label{subsec:sample-proportion-reference}
%%%%%%%%%%%%%%%%%%%%%%%%%%%%%%%%%%%%%%%%%%%%%%%%%%%%%%%%%%%%

If \(x\) observations in a sample of size \(n\) are classified as successes,

\[
\boxed{
\hat p
=
\frac{x}{n}.
}
\]

%%%%%%%%%%%%%%%%%%%%%%%%%%%%%%%%%%%%%%%%%%%%%%%%%%%%%%%%%%%%
\subsection{Covariance}
\label{subsec:sample-covariance-reference}
%%%%%%%%%%%%%%%%%%%%%%%%%%%%%%%%%%%%%%%%%%%%%%%%%%%%%%%%%%%%

Sample covariance is

\[
\boxed{
s_{xy}
=
\frac{
1
}{
n-1
}
\sum_{i=1}^{n}
(x_i-\bar{x})
(y_i-\bar{y}).
}
\]

Population covariance is

\[
\boxed{
\operatorname{Cov}(X,Y)
=
\mathbb{E}
[
(X-\mu_X)(Y-\mu_Y)
].
}
\]

%%%%%%%%%%%%%%%%%%%%%%%%%%%%%%%%%%%%%%%%%%%%%%%%%%%%%%%%%%%%
\subsection{Pearson Correlation}
\label{subsec:pearson-correlation-reference}
%%%%%%%%%%%%%%%%%%%%%%%%%%%%%%%%%%%%%%%%%%%%%%%%%%%%%%%%%%%%

The sample Pearson correlation coefficient is

\[
\boxed{
r
=
\frac{
\sum_i
(x_i-\bar{x})
(y_i-\bar{y})
}{
\sqrt{
\sum_i
(x_i-\bar{x})^2
}
\sqrt{
\sum_i
(y_i-\bar{y})^2
}
}.
}
\]

Equivalently,

\[
\boxed{
r
=
\frac{
s_{xy}
}{
s_xs_y
}.
}
\]

It satisfies

\[
\boxed{
-1\leq r\leq1.
}
\]

%%%%%%%%%%%%%%%%%%%%%%%%%%%%%%%%%%%%%%%%%%%%%%%%%%%%%%%%%%%%
\subsection{Standard Error of the Sample Mean}
\label{subsec:standard-error-mean}
%%%%%%%%%%%%%%%%%%%%%%%%%%%%%%%%%%%%%%%%%%%%%%%%%%%%%%%%%%%%

If population standard deviation \(\sigma\) is known,

\[
\boxed{
\operatorname{SE}(\bar X)
=
\frac{
\sigma
}{
\sqrt n
}.
}
\]

When \(\sigma\) is unknown, the estimated standard error is

\[
\boxed{
\widehat{\operatorname{SE}}(\bar X)
=
\frac{
s
}{
\sqrt n
}.
}
\]

%%%%%%%%%%%%%%%%%%%%%%%%%%%%%%%%%%%%%%%%%%%%%%%%%%%%%%%%%%%%
\subsection{Finite Population Correction}
\label{subsec:finite-population-correction}
%%%%%%%%%%%%%%%%%%%%%%%%%%%%%%%%%%%%%%%%%%%%%%%%%%%%%%%%%%%%

For simple random sampling without replacement from a finite population,

\[
\boxed{
\operatorname{SE}(\bar X)
=
\frac{
\sigma
}{
\sqrt n
}
\sqrt{
\frac{
N-n
}{
N-1
}
}.
}
\]

The factor

\[
\boxed{
\sqrt{
\frac{
N-n
}{
N-1
}
}
}
\]

is the finite-population correction.

%%%%%%%%%%%%%%%%%%%%%%%%%%%%%%%%%%%%%%%%%%%%%%%%%%%%%%%%%%%%
\subsection{Standard Error of a Sample Proportion}
\label{subsec:standard-error-proportion}
%%%%%%%%%%%%%%%%%%%%%%%%%%%%%%%%%%%%%%%%%%%%%%%%%%%%%%%%%%%%

For population proportion \(p\),

\[
\boxed{
\operatorname{SE}(\hat p)
=
\sqrt{
\frac{
p(1-p)
}{
n
}
}.
}
\]

Estimated standard error:

\[
\boxed{
\widehat{\operatorname{SE}}(\hat p)
=
\sqrt{
\frac{
\hat p(1-\hat p)
}{
n
}
}.
}
\]

%%%%%%%%%%%%%%%%%%%%%%%%%%%%%%%%%%%%%%%%%%%%%%%%%%%%%%%%%%%%
\subsection{Standard Normal Critical Values}
\label{subsec:standard-normal-critical-values}
%%%%%%%%%%%%%%%%%%%%%%%%%%%%%%%%%%%%%%%%%%%%%%%%%%%%%%%%%%%%

Frequently used two-sided confidence critical values are approximately:

\[
\boxed{
\begin{array}{c|c}
\text{Confidence Level}
&
z_{\alpha/2}
\\ \hline
80\% & 1.282\\
90\% & 1.645\\
95\% & 1.960\\
98\% & 2.326\\
99\% & 2.576
\end{array}
}
\]

Here \(z_{\alpha/2}\) satisfies

\[
\boxed{
\Pr(Z>z_{\alpha/2})
=
\frac{\alpha}{2}.
}
\]

%%%%%%%%%%%%%%%%%%%%%%%%%%%%%%%%%%%%%%%%%%%%%%%%%%%%%%%%%%%%
\subsection{Confidence Interval General Form}
\label{subsec:confidence-interval-general-form}
%%%%%%%%%%%%%%%%%%%%%%%%%%%%%%%%%%%%%%%%%%%%%%%%%%%%%%%%%%%%

Many confidence intervals have structure

\[
\boxed{
\text{estimate}
\pm
(\text{critical value})
(\text{standard error}).
}
\]

The quantity

\[
\boxed{
(\text{critical value})
(\text{standard error})
}
\]

is the margin of error.

%%%%%%%%%%%%%%%%%%%%%%%%%%%%%%%%%%%%%%%%%%%%%%%%%%%%%%%%%%%%
\subsection{Confidence Interval for a Mean With Known \(\sigma\)}
\label{subsec:ci-mean-known-sigma}
%%%%%%%%%%%%%%%%%%%%%%%%%%%%%%%%%%%%%%%%%%%%%%%%%%%%%%%%%%%%

Under appropriate sampling assumptions,

\[
\boxed{
\bar{x}
\pm
z_{\alpha/2}
\frac{
\sigma
}{
\sqrt n
}.
}
\]

%%%%%%%%%%%%%%%%%%%%%%%%%%%%%%%%%%%%%%%%%%%%%%%%%%%%%%%%%%%%
\subsection{Confidence Interval for a Mean With Unknown \(\sigma\)}
\label{subsec:ci-mean-unknown-sigma}
%%%%%%%%%%%%%%%%%%%%%%%%%%%%%%%%%%%%%%%%%%%%%%%%%%%%%%%%%%%%

For normally distributed data, or under suitable large-sample conditions,

\[
\boxed{
\bar{x}
\pm
t_{\alpha/2,n-1}
\frac{
s
}{
\sqrt n
}.
}
\]

Degrees of freedom:

\[
\boxed{
\nu=n-1.
}
\]

%%%%%%%%%%%%%%%%%%%%%%%%%%%%%%%%%%%%%%%%%%%%%%%%%%%%%%%%%%%%
\subsection{Confidence Interval for a Proportion}
\label{subsec:ci-proportion-reference}
%%%%%%%%%%%%%%%%%%%%%%%%%%%%%%%%%%%%%%%%%%%%%%%%%%%%%%%%%%%%

The common Wald interval is

\[
\boxed{
\hat p
\pm
z_{\alpha/2}
\sqrt{
\frac{
\hat p(1-\hat p)
}{
n
}
}.
}
\]

For small samples or extreme proportions, alternatives such as Wilson or
exact intervals may be preferable.

%%%%%%%%%%%%%%%%%%%%%%%%%%%%%%%%%%%%%%%%%%%%%%%%%%%%%%%%%%%%
\subsection{Difference of Two Means: Known Variances}
\label{subsec:ci-two-means-known-variances}
%%%%%%%%%%%%%%%%%%%%%%%%%%%%%%%%%%%%%%%%%%%%%%%%%%%%%%%%%%%%

For independent samples,

\[
\boxed{
(\bar{x}_1-\bar{x}_2)
\pm
z_{\alpha/2}
\sqrt{
\frac{
\sigma_1^2
}{
n_1
}
+
\frac{
\sigma_2^2
}{
n_2
}
}.
}
\]

%%%%%%%%%%%%%%%%%%%%%%%%%%%%%%%%%%%%%%%%%%%%%%%%%%%%%%%%%%%%
\subsection{Welch Two-Sample \(t\) Interval}
\label{subsec:welch-two-sample-interval}
%%%%%%%%%%%%%%%%%%%%%%%%%%%%%%%%%%%%%%%%%%%%%%%%%%%%%%%%%%%%

For independent samples with unknown, not necessarily equal variances,

\[
\boxed{
(\bar{x}_1-\bar{x}_2)
\pm
t_{\alpha/2,\nu}
\sqrt{
\frac{
s_1^2
}{
n_1
}
+
\frac{
s_2^2
}{
n_2
}
}.
}
\]

Welch degrees of freedom are approximated by

\[
\boxed{
\nu
=
\frac{
\left(
\frac{s_1^2}{n_1}
+
\frac{s_2^2}{n_2}
\right)^2
}{
\frac{
(s_1^2/n_1)^2
}{
n_1-1
}
+
\frac{
(s_2^2/n_2)^2
}{
n_2-1
}
}.
}
\]

%%%%%%%%%%%%%%%%%%%%%%%%%%%%%%%%%%%%%%%%%%%%%%%%%%%%%%%%%%%%
\subsection{Paired-Sample Confidence Interval}
\label{subsec:paired-confidence-interval}
%%%%%%%%%%%%%%%%%%%%%%%%%%%%%%%%%%%%%%%%%%%%%%%%%%%%%%%%%%%%

Let

\[
d_i
=
x_i-y_i.
\]

Then analyze the differences as one sample.

The interval is

\[
\boxed{
\bar d
\pm
t_{\alpha/2,n-1}
\frac{
s_d
}{
\sqrt n
}.
}
\]

%%%%%%%%%%%%%%%%%%%%%%%%%%%%%%%%%%%%%%%%%%%%%%%%%%%%%%%%%%%%
\subsection{One-Sample \(z\)-Test for a Mean}
\label{subsec:one-sample-z-test-mean}
%%%%%%%%%%%%%%%%%%%%%%%%%%%%%%%%%%%%%%%%%%%%%%%%%%%%%%%%%%%%

If \(\sigma\) is known,

\[
\boxed{
z
=
\frac{
\bar{x}-\mu_0
}{
\sigma/\sqrt n
}.
}
\]

%%%%%%%%%%%%%%%%%%%%%%%%%%%%%%%%%%%%%%%%%%%%%%%%%%%%%%%%%%%%
\subsection{One-Sample \(t\)-Test for a Mean}
\label{subsec:one-sample-t-test-mean}
%%%%%%%%%%%%%%%%%%%%%%%%%%%%%%%%%%%%%%%%%%%%%%%%%%%%%%%%%%%%

If \(\sigma\) is unknown,

\[
\boxed{
t
=
\frac{
\bar{x}-\mu_0
}{
s/\sqrt n
}.
}
\]

Degrees of freedom:

\[
\boxed{
\nu=n-1.
}
\]

%%%%%%%%%%%%%%%%%%%%%%%%%%%%%%%%%%%%%%%%%%%%%%%%%%%%%%%%%%%%
\subsection{One-Sample Proportion Test}
\label{subsec:one-sample-proportion-test}
%%%%%%%%%%%%%%%%%%%%%%%%%%%%%%%%%%%%%%%%%%%%%%%%%%%%%%%%%%%%

For null hypothesis

\[
H_0:p=p_0,
\]

the common large-sample test statistic is

\[
\boxed{
z
=
\frac{
\hat p-p_0
}{
\sqrt{
p_0(1-p_0)/n
}
}.
}
\]

The null proportion \(p_0\), not \(\hat p\), appears in the standard error
for this test.

%%%%%%%%%%%%%%%%%%%%%%%%%%%%%%%%%%%%%%%%%%%%%%%%%%%%%%%%%%%%
\subsection{Two-Proportion Test}
\label{subsec:two-proportion-test-reference}
%%%%%%%%%%%%%%%%%%%%%%%%%%%%%%%%%%%%%%%%%%%%%%%%%%%%%%%%%%%%

For testing

\[
H_0:p_1=p_2,
\]

use pooled proportion

\[
\boxed{
\hat p
=
\frac{
x_1+x_2
}{
n_1+n_2
}.
}
\]

Then

\[
\boxed{
z
=
\frac{
\hat p_1-\hat p_2
}{
\sqrt{
\hat p(1-\hat p)
\left(
\frac1{n_1}
+
\frac1{n_2}
\right)
}
}.
}
\]

%%%%%%%%%%%%%%%%%%%%%%%%%%%%%%%%%%%%%%%%%%%%%%%%%%%%%%%%%%%%
\subsection{Hypothesis-Test Structure}
\label{subsec:hypothesis-test-structure}
%%%%%%%%%%%%%%%%%%%%%%%%%%%%%%%%%%%%%%%%%%%%%%%%%%%%%%%%%%%%

A hypothesis test typically contains:

\[
\boxed{
H_0
=
\text{null hypothesis},
}
\]

\[
\boxed{
H_A
=
\text{alternative hypothesis},
}
\]

\[
\boxed{
\alpha
=
\text{significance level},
}
\]

\[
\boxed{
\text{test statistic},
}
\]

and

\[
\boxed{
p\text{-value}.
}
\]

%%%%%%%%%%%%%%%%%%%%%%%%%%%%%%%%%%%%%%%%%%%%%%%%%%%%%%%%%%%%
\subsection{\(p\)-Value}
\label{subsec:p-value-reference}
%%%%%%%%%%%%%%%%%%%%%%%%%%%%%%%%%%%%%%%%%%%%%%%%%%%%%%%%%%%%

The \(p\)-value is the probability, computed under the null model, of
obtaining a test statistic at least as incompatible with \(H_0\) as the
observed statistic, using the chosen test's ordering.

It is not the probability that \(H_0\) is true.

%%%%%%%%%%%%%%%%%%%%%%%%%%%%%%%%%%%%%%%%%%%%%%%%%%%%%%%%%%%%
\subsection{Decision Rule}
\label{subsec:hypothesis-test-decision-rule}
%%%%%%%%%%%%%%%%%%%%%%%%%%%%%%%%%%%%%%%%%%%%%%%%%%%%%%%%%%%%

A common decision rule is:

\[
\boxed{
p\leq\alpha
\Longrightarrow
\text{reject }H_0.
}
\]

If

\[
p>\alpha,
\]

one fails to reject \(H_0\).

This does not prove that \(H_0\) is true.

%%%%%%%%%%%%%%%%%%%%%%%%%%%%%%%%%%%%%%%%%%%%%%%%%%%%%%%%%%%%
\subsection{Type I Error}
\label{subsec:type-one-error-reference}
%%%%%%%%%%%%%%%%%%%%%%%%%%%%%%%%%%%%%%%%%%%%%%%%%%%%%%%%%%%%

A Type I error occurs when \(H_0\) is rejected even though \(H_0\) is true.

Its probability under the designed test is controlled by

\[
\boxed{
\alpha.
}
\]

%%%%%%%%%%%%%%%%%%%%%%%%%%%%%%%%%%%%%%%%%%%%%%%%%%%%%%%%%%%%
\subsection{Type II Error}
\label{subsec:type-two-error-reference}
%%%%%%%%%%%%%%%%%%%%%%%%%%%%%%%%%%%%%%%%%%%%%%%%%%%%%%%%%%%%

A Type II error occurs when \(H_0\) is not rejected even though a specified
alternative is true.

Its probability is commonly denoted

\[
\boxed{
\beta.
}
\]

%%%%%%%%%%%%%%%%%%%%%%%%%%%%%%%%%%%%%%%%%%%%%%%%%%%%%%%%%%%%
\subsection{Statistical Power}
\label{subsec:statistical-power-reference}
%%%%%%%%%%%%%%%%%%%%%%%%%%%%%%%%%%%%%%%%%%%%%%%%%%%%%%%%%%%%

Power is

\[
\boxed{
1-\beta.
}
\]

It is the probability of rejecting \(H_0\) under a specified alternative
model.

%%%%%%%%%%%%%%%%%%%%%%%%%%%%%%%%%%%%%%%%%%%%%%%%%%%%%%%%%%%%
\subsection{Chi-Square Test Statistic}
\label{subsec:chi-square-test-statistic}
%%%%%%%%%%%%%%%%%%%%%%%%%%%%%%%%%%%%%%%%%%%%%%%%%%%%%%%%%%%%

For observed counts \(O_i\) and expected counts \(E_i\),

\[
\boxed{
\chi^2
=
\sum_i
\frac{
(O_i-E_i)^2
}{
E_i
}.
}
\]

This structure appears in goodness-of-fit and contingency-table tests.

%%%%%%%%%%%%%%%%%%%%%%%%%%%%%%%%%%%%%%%%%%%%%%%%%%%%%%%%%%%%
\subsection{Expected Counts in a Contingency Table}
\label{subsec:contingency-expected-count}
%%%%%%%%%%%%%%%%%%%%%%%%%%%%%%%%%%%%%%%%%%%%%%%%%%%%%%%%%%%%

Under independence,

\[
\boxed{
E_{ij}
=
\frac{
(\text{row }i\text{ total})
(\text{column }j\text{ total})
}{
\text{grand total}
}.
}
\]

%%%%%%%%%%%%%%%%%%%%%%%%%%%%%%%%%%%%%%%%%%%%%%%%%%%%%%%%%%%%
\subsection{Degrees of Freedom for an \(r\times c\) Table}
\label{subsec:contingency-degrees-freedom}
%%%%%%%%%%%%%%%%%%%%%%%%%%%%%%%%%%%%%%%%%%%%%%%%%%%%%%%%%%%%

For a standard chi-square test of independence,

\[
\boxed{
\nu
=
(r-1)(c-1).
}
\]

%%%%%%%%%%%%%%%%%%%%%%%%%%%%%%%%%%%%%%%%%%%%%%%%%%%%%%%%%%%%
\subsection{Simple Linear Regression Model}
\label{subsec:simple-linear-regression-reference}
%%%%%%%%%%%%%%%%%%%%%%%%%%%%%%%%%%%%%%%%%%%%%%%%%%%%%%%%%%%%

The standard model is

\[
\boxed{
Y_i
=
\beta_0
+
\beta_1x_i
+
\varepsilon_i.
}
\]

The fitted line is

\[
\boxed{
\hat y
=
b_0+b_1x.
}
\]

%%%%%%%%%%%%%%%%%%%%%%%%%%%%%%%%%%%%%%%%%%%%%%%%%%%%%%%%%%%%
\subsection{Regression Slope Estimate}
\label{subsec:regression-slope-estimate}
%%%%%%%%%%%%%%%%%%%%%%%%%%%%%%%%%%%%%%%%%%%%%%%%%%%%%%%%%%%%

\[
\boxed{
b_1
=
\frac{
\sum_i
(x_i-\bar{x})
(y_i-\bar{y})
}{
\sum_i
(x_i-\bar{x})^2
}.
}
\]

Equivalently,

\[
\boxed{
b_1
=
r
\frac{
s_y
}{
s_x
}.
}
\]

%%%%%%%%%%%%%%%%%%%%%%%%%%%%%%%%%%%%%%%%%%%%%%%%%%%%%%%%%%%%
\subsection{Regression Intercept Estimate}
\label{subsec:regression-intercept-estimate}
%%%%%%%%%%%%%%%%%%%%%%%%%%%%%%%%%%%%%%%%%%%%%%%%%%%%%%%%%%%%

\[
\boxed{
b_0
=
\bar{y}
-
b_1\bar{x}.
}
\]

Thus the least-squares regression line passes through

\[
\boxed{
(\bar{x},\bar{y}).
}
\]

%%%%%%%%%%%%%%%%%%%%%%%%%%%%%%%%%%%%%%%%%%%%%%%%%%%%%%%%%%%%
\subsection{Residuals}
\label{subsec:residual-reference}
%%%%%%%%%%%%%%%%%%%%%%%%%%%%%%%%%%%%%%%%%%%%%%%%%%%%%%%%%%%%

The residual for observation \(i\) is

\[
\boxed{
e_i
=
y_i-\hat y_i.
}
\]

%%%%%%%%%%%%%%%%%%%%%%%%%%%%%%%%%%%%%%%%%%%%%%%%%%%%%%%%%%%%
\subsection{Regression Sums of Squares}
\label{subsec:regression-sums-squares}
%%%%%%%%%%%%%%%%%%%%%%%%%%%%%%%%%%%%%%%%%%%%%%%%%%%%%%%%%%%%

Total sum of squares:

\[
\boxed{
SST
=
\sum_i
(y_i-\bar y)^2.
}
\]

Regression sum of squares:

\[
\boxed{
SSR
=
\sum_i
(\hat y_i-\bar y)^2.
}
\]

Error sum of squares:

\[
\boxed{
SSE
=
\sum_i
(y_i-\hat y_i)^2.
}
\]

For ordinary least squares with an intercept,

\[
\boxed{
SST
=
SSR+SSE.
}
\]

%%%%%%%%%%%%%%%%%%%%%%%%%%%%%%%%%%%%%%%%%%%%%%%%%%%%%%%%%%%%
\subsection{Coefficient of Determination}
\label{subsec:r-squared-reference}
%%%%%%%%%%%%%%%%%%%%%%%%%%%%%%%%%%%%%%%%%%%%%%%%%%%%%%%%%%%%

\[
\boxed{
R^2
=
\frac{
SSR
}{
SST
}
=
1-
\frac{
SSE
}{
SST
}.
}
\]

For simple linear regression with an intercept,

\[
\boxed{
R^2=r^2.
}
\]

%%%%%%%%%%%%%%%%%%%%%%%%%%%%%%%%%%%%%%%%%%%%%%%%%%%%%%%%%%%%
\subsection{Residual Standard Error}
\label{subsec:residual-standard-error}
%%%%%%%%%%%%%%%%%%%%%%%%%%%%%%%%%%%%%%%%%%%%%%%%%%%%%%%%%%%%

For simple linear regression,

\[
\boxed{
s_e
=
\sqrt{
\frac{
SSE
}{
n-2
}
}.
}
\]

%%%%%%%%%%%%%%%%%%%%%%%%%%%%%%%%%%%%%%%%%%%%%%%%%%%%%%%%%%%%
\subsection{Standard Error of Regression Slope}
\label{subsec:slope-standard-error}
%%%%%%%%%%%%%%%%%%%%%%%%%%%%%%%%%%%%%%%%%%%%%%%%%%%%%%%%%%%%

\[
\boxed{
\operatorname{SE}(b_1)
=
\frac{
s_e
}{
\sqrt{
\sum_i
(x_i-\bar{x})^2
}
}.
}
\]

A test statistic for

\[
H_0:\beta_1=0
\]

is

\[
\boxed{
t
=
\frac{
b_1
}{
\operatorname{SE}(b_1)
},
}
\]

with

\[
\boxed{
n-2
}
\]

degrees of freedom.

%%%%%%%%%%%%%%%%%%%%%%%%%%%%%%%%%%%%%%%%%%%%%%%%%%%%%%%%%%%%
\subsection{ANOVA Decomposition}
\label{subsec:anova-decomposition-reference}
%%%%%%%%%%%%%%%%%%%%%%%%%%%%%%%%%%%%%%%%%%%%%%%%%%%%%%%%%%%%

In one-way ANOVA,

\[
\boxed{
SS_T
=
SS_B
+
SS_W,
}
\]

where

\[
SS_T
=
\text{total sum of squares},
\]

\[
SS_B
=
\text{between-group sum of squares},
\]

and

\[
SS_W
=
\text{within-group sum of squares}.
\]

%%%%%%%%%%%%%%%%%%%%%%%%%%%%%%%%%%%%%%%%%%%%%%%%%%%%%%%%%%%%
\subsection{One-Way ANOVA Between-Groups Sum of Squares}
\label{subsec:anova-between-groups}
%%%%%%%%%%%%%%%%%%%%%%%%%%%%%%%%%%%%%%%%%%%%%%%%%%%%%%%%%%%%

For group means \(\bar{x}_j\), group sizes \(n_j\), and grand mean
\(\bar{x}\),

\[
\boxed{
SS_B
=
\sum_{j=1}^{k}
n_j
(\bar{x}_j-\bar{x})^2.
}
\]

%%%%%%%%%%%%%%%%%%%%%%%%%%%%%%%%%%%%%%%%%%%%%%%%%%%%%%%%%%%%
\subsection{One-Way ANOVA Within-Groups Sum of Squares}
\label{subsec:anova-within-groups}
%%%%%%%%%%%%%%%%%%%%%%%%%%%%%%%%%%%%%%%%%%%%%%%%%%%%%%%%%%%%

\[
\boxed{
SS_W
=
\sum_{j=1}^{k}
\sum_{i=1}^{n_j}
(x_{ij}-\bar{x}_j)^2.
}
\]

%%%%%%%%%%%%%%%%%%%%%%%%%%%%%%%%%%%%%%%%%%%%%%%%%%%%%%%%%%%%
\subsection{ANOVA Degrees of Freedom}
\label{subsec:anova-degrees-freedom}
%%%%%%%%%%%%%%%%%%%%%%%%%%%%%%%%%%%%%%%%%%%%%%%%%%%%%%%%%%%%

Between groups:

\[
\boxed{
df_B
=
k-1.
}
\]

Within groups:

\[
\boxed{
df_W
=
N-k.
}
\]

Total:

\[
\boxed{
df_T
=
N-1.
}
\]

%%%%%%%%%%%%%%%%%%%%%%%%%%%%%%%%%%%%%%%%%%%%%%%%%%%%%%%%%%%%
\subsection{ANOVA Mean Squares}
\label{subsec:anova-mean-squares}
%%%%%%%%%%%%%%%%%%%%%%%%%%%%%%%%%%%%%%%%%%%%%%%%%%%%%%%%%%%%

\[
\boxed{
MS_B
=
\frac{
SS_B
}{
k-1
}.
}
\]

\[
\boxed{
MS_W
=
\frac{
SS_W
}{
N-k
}.
}
\]

%%%%%%%%%%%%%%%%%%%%%%%%%%%%%%%%%%%%%%%%%%%%%%%%%%%%%%%%%%%%
\subsection{ANOVA \(F\)-Statistic}
\label{subsec:anova-f-statistic}
%%%%%%%%%%%%%%%%%%%%%%%%%%%%%%%%%%%%%%%%%%%%%%%%%%%%%%%%%%%%

\[
\boxed{
F
=
\frac{
MS_B
}{
MS_W
}.
}
\]

Under the standard null model and assumptions, this is compared with an
\(F\) distribution.

%%%%%%%%%%%%%%%%%%%%%%%%%%%%%%%%%%%%%%%%%%%%%%%%%%%%%%%%%%%%
\subsection{Odds}
\label{subsec:odds-reference}
%%%%%%%%%%%%%%%%%%%%%%%%%%%%%%%%%%%%%%%%%%%%%%%%%%%%%%%%%%%%

For event probability \(p\),

\[
\boxed{
\text{odds}
=
\frac{
p
}{
1-p
}.
}
\]

Conversely,

\[
\boxed{
p
=
\frac{
\text{odds}
}{
1+\text{odds}
}.
}
\]

%%%%%%%%%%%%%%%%%%%%%%%%%%%%%%%%%%%%%%%%%%%%%%%%%%%%%%%%%%%%
\subsection{Odds Ratio}
\label{subsec:odds-ratio-reference}
%%%%%%%%%%%%%%%%%%%%%%%%%%%%%%%%%%%%%%%%%%%%%%%%%%%%%%%%%%%%

For two groups with probabilities \(p_1,p_2\),

\[
\boxed{
OR
=
\frac{
p_1/(1-p_1)
}{
p_2/(1-p_2)
}.
}
\]

%%%%%%%%%%%%%%%%%%%%%%%%%%%%%%%%%%%%%%%%%%%%%%%%%%%%%%%%%%%%
\subsection{Relative Risk}
\label{subsec:relative-risk-reference}
%%%%%%%%%%%%%%%%%%%%%%%%%%%%%%%%%%%%%%%%%%%%%%%%%%%%%%%%%%%%

For risks \(p_1,p_2\),

\[
\boxed{
RR
=
\frac{
p_1
}{
p_2
}.
}
\]

Relative risk and odds ratio are not the same quantity.

%%%%%%%%%%%%%%%%%%%%%%%%%%%%%%%%%%%%%%%%%%%%%%%%%%%%%%%%%%%%
\subsection{Risk Difference}
\label{subsec:risk-difference-reference}
%%%%%%%%%%%%%%%%%%%%%%%%%%%%%%%%%%%%%%%%%%%%%%%%%%%%%%%%%%%%

\[
\boxed{
RD
=
p_1-p_2.
}
\]

This is an absolute effect measure.

%%%%%%%%%%%%%%%%%%%%%%%%%%%%%%%%%%%%%%%%%%%%%%%%%%%%%%%%%%%%
\subsection{Cohen's \(d\)}
\label{subsec:cohens-d-reference}
%%%%%%%%%%%%%%%%%%%%%%%%%%%%%%%%%%%%%%%%%%%%%%%%%%%%%%%%%%%%

For two groups, a common standardized mean difference is

\[
\boxed{
d
=
\frac{
\bar{x}_1-\bar{x}_2
}{
s_{\mathrm{pooled}}
}.
}
\]

For equal-variance settings,

\[
\boxed{
s_{\mathrm{pooled}}
=
\sqrt{
\frac{
(n_1-1)s_1^2
+
(n_2-1)s_2^2
}{
n_1+n_2-2
}
}.
}
\]

%%%%%%%%%%%%%%%%%%%%%%%%%%%%%%%%%%%%%%%%%%%%%%%%%%%%%%%%%%%%
\subsection{Eta Squared}
\label{subsec:eta-squared-reference}
%%%%%%%%%%%%%%%%%%%%%%%%%%%%%%%%%%%%%%%%%%%%%%%%%%%%%%%%%%%%

A common ANOVA effect-size measure is

\[
\boxed{
\eta^2
=
\frac{
SS_B
}{
SS_T
}.
}
\]

%%%%%%%%%%%%%%%%%%%%%%%%%%%%%%%%%%%%%%%%%%%%%%%%%%%%%%%%%%%%
\subsection{Coefficient of Variation}
\label{subsec:coefficient-variation-reference}
%%%%%%%%%%%%%%%%%%%%%%%%%%%%%%%%%%%%%%%%%%%%%%%%%%%%%%%%%%%%

For positive-scale data with meaningful zero,

\[
\boxed{
CV
=
\frac{
s
}{
\bar{x}
}
}
\]

or as a percentage,

\[
\boxed{
100
\frac{
s
}{
\bar{x}
}
\%.
}
\]

The coefficient of variation is not appropriate for every measurement scale.

%%%%%%%%%%%%%%%%%%%%%%%%%%%%%%%%%%%%%%%%%%%%%%%%%%%%%%%%%%%%
\subsection{Standard Error Versus Standard Deviation}
\label{subsec:se-versus-sd}
%%%%%%%%%%%%%%%%%%%%%%%%%%%%%%%%%%%%%%%%%%%%%%%%%%%%%%%%%%%%

Standard deviation measures variability among observations.

Standard error measures uncertainty or sampling variability in an estimator.

For a sample mean,

\[
\boxed{
SE
\propto
\frac1{\sqrt n}.
}
\]

Thus standard error generally decreases as sample size increases.

%%%%%%%%%%%%%%%%%%%%%%%%%%%%%%%%%%%%%%%%%%%%%%%%%%%%%%%%%%%%
\subsection{Critical Values and Quantiles}
\label{subsec:critical-values-quantiles}
%%%%%%%%%%%%%%%%%%%%%%%%%%%%%%%%%%%%%%%%%%%%%%%%%%%%%%%%%%%%

For a distribution with CDF \(F\), a quantile \(q_p\) satisfies

\[
\boxed{
F(q_p)=p
}
\]

when the inverse is uniquely defined.

Statistical critical values are distribution quantiles chosen according to a
significance level.

%%%%%%%%%%%%%%%%%%%%%%%%%%%%%%%%%%%%%%%%%%%%%%%%%%%%%%%%%%%%
\subsection{Common \(t\)-Critical Values for 95\% Confidence}
\label{subsec:t-critical-selected-values}
%%%%%%%%%%%%%%%%%%%%%%%%%%%%%%%%%%%%%%%%%%%%%%%%%%%%%%%%%%%%

Selected approximate values of

\[
t_{0.025,\nu}
\]

are:

\[
\boxed{
\begin{array}{c|c}
\nu & t_{0.025,\nu} \\ \hline
1 & 12.706\\
2 & 4.303\\
5 & 2.571\\
10 & 2.228\\
20 & 2.086\\
30 & 2.042\\
60 & 2.000\\
120 & 1.980\\
\infty & 1.960
\end{array}
}
\]

The \(\nu=\infty\) row corresponds to the standard normal limit.

%%%%%%%%%%%%%%%%%%%%%%%%%%%%%%%%%%%%%%%%%%%%%%%%%%%%%%%%%%%%
\subsection{Bonferroni Correction}
\label{subsec:bonferroni-reference}
%%%%%%%%%%%%%%%%%%%%%%%%%%%%%%%%%%%%%%%%%%%%%%%%%%%%%%%%%%%%

For \(m\) simultaneous tests and desired familywise level \(\alpha\), the
Bonferroni method uses individual significance level

\[
\boxed{
\frac{\alpha}{m}.
}
\]

Equivalently, multiply each raw \(p\)-value by \(m\) and cap at \(1\).

%%%%%%%%%%%%%%%%%%%%%%%%%%%%%%%%%%%%%%%%%%%%%%%%%%%%%%%%%%%%
\subsection{False Discovery Rate}
\label{subsec:false-discovery-rate-reference}
%%%%%%%%%%%%%%%%%%%%%%%%%%%%%%%%%%%%%%%%%%%%%%%%%%%%%%%%%%%%

False discovery rate procedures, such as the Benjamini--Hochberg method, aim
to control the expected proportion of false discoveries among rejected
hypotheses.

They address a different error criterion from familywise error control.

%%%%%%%%%%%%%%%%%%%%%%%%%%%%%%%%%%%%%%%%%%%%%%%%%%%%%%%%%%%%
\subsection{Benjamini--Hochberg Threshold}
\label{subsec:benjamini-hochberg-reference}
%%%%%%%%%%%%%%%%%%%%%%%%%%%%%%%%%%%%%%%%%%%%%%%%%%%%%%%%%%%%

Order the \(m\) \(p\)-values:

\[
p_{(1)}
\leq
p_{(2)}
\leq
\cdots
\leq
p_{(m)}.
\]

For target level \(q\), find the largest \(k\) satisfying

\[
\boxed{
p_{(k)}
\leq
\frac{k}{m}q.
}
\]

Reject the hypotheses corresponding to

\[
p_{(1)},\ldots,p_{(k)}
\]

under the standard Benjamini--Hochberg procedure and its usual conditions.

%%%%%%%%%%%%%%%%%%%%%%%%%%%%%%%%%%%%%%%%%%%%%%%%%%%%%%%%%%%%
\subsection{Quick Descriptive Statistics Table}
\label{subsec:quick-descriptive-statistics-table}
%%%%%%%%%%%%%%%%%%%%%%%%%%%%%%%%%%%%%%%%%%%%%%%%%%%%%%%%%%%%

\[
\boxed{
\bar{x}
=
\frac1n
\sum x_i
}
\]

\[
\boxed{
s^2
=
\frac{
\sum
(x_i-\bar{x})^2
}{
n-1
}
}
\]

\[
\boxed{
s
=
\sqrt{s^2}
}
\]

\[
\boxed{
\operatorname{IQR}
=
Q_3-Q_1
}
\]

\[
\boxed{
z
=
\frac{x-\mu}{\sigma}
}
\]

\[
\boxed{
\hat p
=
\frac{x}{n}.
}
\]

%%%%%%%%%%%%%%%%%%%%%%%%%%%%%%%%%%%%%%%%%%%%%%%%%%%%%%%%%%%%
\subsection{Quick Standard Error Table}
\label{subsec:quick-standard-error-table}
%%%%%%%%%%%%%%%%%%%%%%%%%%%%%%%%%%%%%%%%%%%%%%%%%%%%%%%%%%%%

Mean, known \(\sigma\):

\[
\boxed{
SE(\bar X)
=
\frac{\sigma}{\sqrt n}.
}
\]

Mean, estimated:

\[
\boxed{
\widehat{SE}(\bar X)
=
\frac{s}{\sqrt n}.
}
\]

Proportion:

\[
\boxed{
\widehat{SE}(\hat p)
=
\sqrt{
\frac{
\hat p(1-\hat p)
}{
n
}
}.
}
\]

Difference of independent means:

\[
\boxed{
\widehat{SE}
(\bar X_1-\bar X_2)
=
\sqrt{
\frac{s_1^2}{n_1}
+
\frac{s_2^2}{n_2}
}.
}
\]

%%%%%%%%%%%%%%%%%%%%%%%%%%%%%%%%%%%%%%%%%%%%%%%%%%%%%%%%%%%%
\subsection{Quick Confidence Interval Table}
\label{subsec:quick-confidence-interval-table}
%%%%%%%%%%%%%%%%%%%%%%%%%%%%%%%%%%%%%%%%%%%%%%%%%%%%%%%%%%%%

Mean with known \(\sigma\):

\[
\boxed{
\bar{x}
\pm
z_{\alpha/2}
\frac{\sigma}{\sqrt n}
}
\]

Mean with unknown \(\sigma\):

\[
\boxed{
\bar{x}
\pm
t_{\alpha/2,n-1}
\frac{s}{\sqrt n}
}
\]

Proportion:

\[
\boxed{
\hat p
\pm
z_{\alpha/2}
\sqrt{
\frac{
\hat p(1-\hat p)
}{
n
}
}
}
\]

Paired mean difference:

\[
\boxed{
\bar d
\pm
t_{\alpha/2,n-1}
\frac{s_d}{\sqrt n}.
}
\]

%%%%%%%%%%%%%%%%%%%%%%%%%%%%%%%%%%%%%%%%%%%%%%%%%%%%%%%%%%%%
\subsection{Quick Test Statistic Table}
\label{subsec:quick-test-statistic-table}
%%%%%%%%%%%%%%%%%%%%%%%%%%%%%%%%%%%%%%%%%%%%%%%%%%%%%%%%%%%%

One mean, known \(\sigma\):

\[
\boxed{
z
=
\frac{
\bar{x}-\mu_0
}{
\sigma/\sqrt n
}
}
\]

One mean, unknown \(\sigma\):

\[
\boxed{
t
=
\frac{
\bar{x}-\mu_0
}{
s/\sqrt n
}
}
\]

One proportion:

\[
\boxed{
z
=
\frac{
\hat p-p_0
}{
\sqrt{
p_0(1-p_0)/n
}
}
}
\]

Chi-square:

\[
\boxed{
\chi^2
=
\sum
\frac{
(O-E)^2
}{
E
}
}
\]

ANOVA:

\[
\boxed{
F
=
\frac{
MS_B
}{
MS_W
}.
}
\]

%%%%%%%%%%%%%%%%%%%%%%%%%%%%%%%%%%%%%%%%%%%%%%%%%%%%%%%%%%%%
\subsection{Worked Example: One-Sample Confidence Interval}
\label{subsec:worked-one-sample-ci}
%%%%%%%%%%%%%%%%%%%%%%%%%%%%%%%%%%%%%%%%%%%%%%%%%%%%%%%%%%%%

\begin{workedexamplebox}{95\% Confidence Interval for a Mean}

Suppose

\[
n=25,
\qquad
\bar{x}=50,
\qquad
s=10.
\]

Using a one-sample \(t\) interval,

\[
\nu=24.
\]

Let

\[
t_{0.025,24}
\approx
2.064.
\]

The standard error is

\[
\frac{10}{\sqrt{25}}
=
2.
\]

Thus the margin of error is

\[
2.064(2)
=
4.128.
\]

Therefore,

\begin{answerbox}
\[
\boxed{
50\pm4.128
}
\]

or approximately

\[
\boxed{
(45.872,\ 54.128).
}
\]
\end{answerbox}

\end{workedexamplebox}

%%%%%%%%%%%%%%%%%%%%%%%%%%%%%%%%%%%%%%%%%%%%%%%%%%%%%%%%%%%%
\subsection{Worked Example: One-Sample \(t\)-Statistic}
\label{subsec:worked-one-sample-t}
%%%%%%%%%%%%%%%%%%%%%%%%%%%%%%%%%%%%%%%%%%%%%%%%%%%%%%%%%%%%

\begin{workedexamplebox}{Test a Mean}

Suppose

\[
\bar{x}=105,
\qquad
\mu_0=100,
\qquad
s=12,
\qquad
n=36.
\]

Then

\[
t
=
\frac{
105-100
}{
12/\sqrt{36}
}.
\]

Since

\[
12/\sqrt{36}=2,
\]

we obtain

\begin{answerbox}
\[
\boxed{
t=2.5.
}
\]
\end{answerbox}

\end{workedexamplebox}

%%%%%%%%%%%%%%%%%%%%%%%%%%%%%%%%%%%%%%%%%%%%%%%%%%%%%%%%%%%%
\subsection{Worked Example: Regression Slope}
\label{subsec:worked-regression-slope}
%%%%%%%%%%%%%%%%%%%%%%%%%%%%%%%%%%%%%%%%%%%%%%%%%%%%%%%%%%%%

\begin{workedexamplebox}{Use Correlation and Standard Deviations}

Suppose

\[
r=0.8,
\qquad
s_x=2,
\qquad
s_y=5.
\]

Then

\[
b_1
=
r
\frac{s_y}{s_x}.
\]

Therefore,

\begin{answerbox}
\[
\boxed{
b_1
=
0.8
\frac52
=
2.
}
\]
\end{answerbox}

\end{workedexamplebox}

%%%%%%%%%%%%%%%%%%%%%%%%%%%%%%%%%%%%%%%%%%%%%%%%%%%%%%%%%%%%
\subsection{Common Errors}
\label{subsec:statistical-tables-common-errors}
%%%%%%%%%%%%%%%%%%%%%%%%%%%%%%%%%%%%%%%%%%%%%%%%%%%%%%%%%%%%

\subsubsection{Using \(n\) Instead of \(n-1\) for Sample Variance}

The usual sample variance estimator is

\[
\boxed{
s^2
=
\frac{
1
}{
n-1
}
\sum
(x_i-\bar{x})^2.
}
\]

%%%%%%%%%%%%%%%%%%%%%%%%%%%%%%%%%%%%%%%%%%%%%%%%%%%%%%%%%%%%
\subsubsection{Confusing Standard Deviation and Standard Error}

Standard deviation describes data spread.

Standard error describes estimator variability.

They are not interchangeable.

%%%%%%%%%%%%%%%%%%%%%%%%%%%%%%%%%%%%%%%%%%%%%%%%%%%%%%%%%%%%
\subsubsection{Using \(z\) Instead of \(t\) Automatically}

When estimating a mean with unknown population standard deviation, the
Student's \(t\) distribution is commonly appropriate under the usual
assumptions.

%%%%%%%%%%%%%%%%%%%%%%%%%%%%%%%%%%%%%%%%%%%%%%%%%%%%%%%%%%%%
\subsubsection{Using \(\hat p\) in the Null Standard Error for a One-Proportion Test}

For

\[
H_0:p=p_0,
\]

the standard null-test denominator is based on

\[
\boxed{
p_0.
}
\]

%%%%%%%%%%%%%%%%%%%%%%%%%%%%%%%%%%%%%%%%%%%%%%%%%%%%%%%%%%%%
\subsubsection{Interpreting a \(p\)-Value as \(\Pr(H_0\text{ is true})\)}

A \(p\)-value is computed assuming the null model.

It is not the posterior probability of the null hypothesis.

%%%%%%%%%%%%%%%%%%%%%%%%%%%%%%%%%%%%%%%%%%%%%%%%%%%%%%%%%%%%
\subsubsection{Saying ``Accept the Null Hypothesis''}

A nonsignificant test generally leads to

\[
\boxed{
\text{fail to reject }H_0,
}
\]

not proof that \(H_0\) is true.

%%%%%%%%%%%%%%%%%%%%%%%%%%%%%%%%%%%%%%%%%%%%%%%%%%%%%%%%%%%%
\subsubsection{Confusing Statistical and Practical Significance}

A very small effect can be statistically significant with a large sample.

A large practical effect may fail to reach significance with limited data.

Effect size and uncertainty should accompany significance testing.

%%%%%%%%%%%%%%%%%%%%%%%%%%%%%%%%%%%%%%%%%%%%%%%%%%%%%%%%%%%%
\subsubsection{Interpreting Correlation as Causation}

A large

\[
|r|
\]

does not prove a causal relationship.

%%%%%%%%%%%%%%%%%%%%%%%%%%%%%%%%%%%%%%%%%%%%%%%%%%%%%%%%%%%%
\subsubsection{Assuming \(R^2\) Measures Causation}

\[
R^2
\]

describes fitted variation explained within the model structure.

It does not establish a causal mechanism.

%%%%%%%%%%%%%%%%%%%%%%%%%%%%%%%%%%%%%%%%%%%%%%%%%%%%%%%%%%%%
\subsubsection{Ignoring Model Assumptions}

Classical procedures may rely on assumptions involving:

\begin{itemize}
    \item independence;
    \item random sampling or randomization;
    \item distributional form;
    \item equal variances;
    \item linearity;
    \item appropriate expected cell counts.
\end{itemize}

The exact assumptions depend on the procedure.

%%%%%%%%%%%%%%%%%%%%%%%%%%%%%%%%%%%%%%%%%%%%%%%%%%%%%%%%%%%%
\subsubsection{Ignoring Multiple Testing}

Performing many hypothesis tests increases opportunities for false
discoveries.

Multiplicity adjustments may be necessary.

%%%%%%%%%%%%%%%%%%%%%%%%%%%%%%%%%%%%%%%%%%%%%%%%%%%%%%%%%%%%
\subsection{Statistical Workflow}
\label{subsec:statistical-reference-workflow}
%%%%%%%%%%%%%%%%%%%%%%%%%%%%%%%%%%%%%%%%%%%%%%%%%%%%%%%%%%%%

A practical statistical workflow is:

\begin{enumerate}
    \item identify the study design and data structure;
    \item distinguish population parameters from sample statistics;
    \item summarize the data descriptively;
    \item inspect distributions and unusual observations;
    \item identify the estimand or scientific question;
    \item choose a model or inferential procedure;
    \item check assumptions;
    \item calculate the estimate;
    \item calculate the standard error;
    \item report confidence intervals where useful;
    \item conduct hypothesis tests only when they answer the question;
    \item report effect sizes;
    \item interpret practical importance;
    \item avoid causal conclusions unless the design supports them.
\end{enumerate}

%%%%%%%%%%%%%%%%%%%%%%%%%%%%%%%%%%%%%%%%%%%%%%%%%%%%%%%%%%%%
\subsection{Exercises}
\label{subsec:statistical-tables-exercises}
%%%%%%%%%%%%%%%%%%%%%%%%%%%%%%%%%%%%%%%%%%%%%%%%%%%%%%%%%%%%

\begin{exercise}[1]
Compute the mean of

\[
2,4,6,8,10.
\]
\end{exercise}

\begin{exercise}[1]
Compute the sample variance of the same data.
\end{exercise}

\begin{exercise}[1]
Compute the sample standard deviation.
\end{exercise}

\begin{exercise}[1]
Compute the range.
\end{exercise}

\begin{exercise}[1]
Compute the sample proportion if \(18\) of \(30\) observations are successes.
\end{exercise}

\begin{exercise}[2]
Compute the \(z\)-score for

\[
x=85,
\qquad
\mu=70,
\qquad
\sigma=10.
\]
\end{exercise}

\begin{exercise}[2]
Compute a weighted mean for specified observations and weights.
\end{exercise}

\begin{exercise}[2]
Compute sample covariance for paired data.
\end{exercise}

\begin{exercise}[2]
Compute Pearson correlation for paired data.
\end{exercise}

\begin{exercise}[2]
Compute the standard error of a mean with

\[
s=12,
\qquad
n=36.
\]
\end{exercise}

\begin{exercise}[2]
Compute the estimated standard error of

\[
\hat p=0.40,
\qquad
n=100.
\]
\end{exercise}

\begin{exercise}[2]
Construct a 95\% confidence interval for a mean with known \(\sigma\).
\end{exercise}

\begin{exercise}[2]
Construct a 95\% one-sample \(t\) interval.
\end{exercise}

\begin{exercise}[2]
Construct a confidence interval for a population proportion.
\end{exercise}

\begin{exercise}[2]
Construct a paired-sample confidence interval.
\end{exercise}

\begin{exercise}[2]
Compute a one-sample \(z\)-test statistic.
\end{exercise}

\begin{exercise}[2]
Compute a one-sample \(t\)-test statistic.
\end{exercise}

\begin{exercise}[2]
Compute a one-proportion test statistic.
\end{exercise}

\begin{exercise}[2]
Compute a two-proportion pooled estimate.
\end{exercise}

\begin{exercise}[2]
Interpret a \(p\)-value in context.
\end{exercise}

\begin{exercise}[2]
Explain the difference between Type I and Type II errors.
\end{exercise}

\begin{exercise}[2]
Compute power when \(\beta=0.20\).
\end{exercise}

\begin{exercise}[2]
Compute expected counts for a \(2\times3\) contingency table.
\end{exercise}

\begin{exercise}[2]
Compute a chi-square statistic from observed and expected counts.
\end{exercise}

\begin{exercise}[2]
Find the degrees of freedom for a \(4\times5\) contingency table.
\end{exercise}

\begin{exercise}[2]
Compute the slope of a simple regression line from summary statistics.
\end{exercise}

\begin{exercise}[2]
Compute the regression intercept.
\end{exercise}

\begin{exercise}[2]
Compute a residual.
\end{exercise}

\begin{exercise}[2]
Verify

\[
SST=SSR+SSE
\]

for a small regression example.
\end{exercise}

\begin{exercise}[2]
Compute

\[
R^2
\]

from \(SSE\) and \(SST\).
\end{exercise}

\begin{exercise}[2]
Compute ANOVA between-group degrees of freedom.
\end{exercise}

\begin{exercise}[2]
Compute ANOVA within-group degrees of freedom.
\end{exercise}

\begin{exercise}[2]
Compute

\[
F
=
MS_B/MS_W.
\]
\end{exercise}

\begin{exercise}[2]
Convert a probability to odds.
\end{exercise}

\begin{exercise}[2]
Convert odds to probability.
\end{exercise}

\begin{exercise}[2]
Compute a relative risk.
\end{exercise}

\begin{exercise}[2]
Compute an odds ratio.
\end{exercise}

\begin{exercise}[2]
Compute Cohen's \(d\) using a pooled standard deviation.
\end{exercise}

\begin{exercise}[2]
Compute eta squared from ANOVA sums of squares.
\end{exercise}

\begin{exercise}[3]
Derive the computational identity for sample sums of squares.
\end{exercise}

\begin{exercise}[3]
Derive the standard error of the sample mean for independent observations.
\end{exercise}

\begin{exercise}[3]
Derive the standard error of a sample proportion from a Bernoulli model.
\end{exercise}

\begin{exercise}[3]
Derive the finite-population correction for simple random sampling without
replacement.
\end{exercise}

\begin{exercise}[3]
Derive the one-sample \(z\)-test statistic from standardization.
\end{exercise}

\begin{exercise}[3]
Derive the pooled standard error for a two-proportion null test.
\end{exercise}

\begin{exercise}[3]
Show that the ordinary least-squares regression line passes through

\[
(\bar{x},\bar{y}).
\]
\end{exercise}

\begin{exercise}[3]
Prove that in simple linear regression with an intercept,

\[
R^2=r^2.
\]
\end{exercise}

\begin{exercise}[3]
Derive the ANOVA decomposition

\[
SS_T=SS_B+SS_W.
\]
\end{exercise}

\begin{exercise}[3]
Compare Bonferroni familywise-error control with false-discovery-rate control.
\end{exercise}

\begin{exercise}[4]
Derive the sampling distribution of the standardized sample mean under a
normal population with known variance.
\end{exercise}

\begin{exercise}[4]
Derive the Student's \(t\) statistic for a normal population with unknown
variance.
\end{exercise}

\begin{exercise}[4]
Derive the chi-square statistic for a normal-population variance.
\end{exercise}

\begin{exercise}[4]
Study the Welch--Satterthwaite degrees-of-freedom approximation.
\end{exercise}

\begin{exercise}[4]
Study the effect of sample size on confidence-interval width.
\end{exercise}

\begin{exercise}[4]
Study the relationship among effect size, sample size, significance level,
and statistical power.
\end{exercise}

%%%%%%%%%%%%%%%%%%%%%%%%%%%%%%%%%%%%%%%%%%%%%%%%%%%%%%%%%%%%
\subsection{Conceptual Summary}
\label{subsec:statistical-tables-summary}
%%%%%%%%%%%%%%%%%%%%%%%%%%%%%%%%%%%%%%%%%%%%%%%%%%%%%%%%%%%%

Descriptive statistics summarize observed data.

The sample mean is

\[
\boxed{
\bar{x}
=
\frac1n
\sum x_i,
}
\]

and sample variance is

\[
\boxed{
s^2
=
\frac{
\sum(x_i-\bar{x})^2
}{
n-1
}.
}
\]

A standard error measures sampling variability of an estimator.

Many confidence intervals have form

\[
\boxed{
\text{estimate}
\pm
\text{critical value}
\times
\text{standard error}.
}
\]

Common test statistics use

\[
\boxed{
z,
\quad
t,
\quad
\chi^2,
\quad
F.
}
\]

Simple regression uses

\[
\boxed{
\hat y
=
b_0+b_1x,
}
\]

with

\[
\boxed{
b_1
=
\frac{
\sum(x_i-\bar{x})(y_i-\bar{y})
}{
\sum(x_i-\bar{x})^2
}.
}
\]

The ANOVA \(F\)-statistic is

\[
\boxed{
F
=
\frac{
MS_B
}{
MS_W
}.
}
\]

The most important interpretation rules are:

\[
\boxed{
p\text{-value}
\neq
\Pr(H_0\text{ is true}),
}
\]

\[
\boxed{
\text{correlation does not imply causation},
}
\]

and

\[
\boxed{
\text{statistical significance does not automatically imply practical
importance}.
}
\]

%%%%%%%%%%%%%%%%%%%%%%%%%%%%%%%%%%%%%%%%%%%%%%%%%%%%%%%%%%%%
\subsection{Section Connections}
\label{subsec:statistical-tables-connections}
%%%%%%%%%%%%%%%%%%%%%%%%%%%%%%%%%%%%%%%%%%%%%%%%%%%%%%%%%%%%

\chapterconnection{
Statistical Tables consolidates the formulas used throughout Descriptive
Statistics, Estimation, Confidence Intervals, Hypothesis Testing, Regression,
Categorical Data Analysis, Experimental Design, ANOVA, and Statistical
Learning. Probability distributions from
Section~\ref{sec:probability-distributions} provide the sampling models behind
\(z\), \(t\), chi-square, and \(F\) procedures; regression connects statistics
with Linear Algebra and Optimization; and effect-size and uncertainty
measures support practical interpretation beyond significance testing. The
next reference section, Common Mathematical Theorems, collects major results
used throughout Algebra, Analysis, Geometry, Topology, Probability, and
Applied Mathematics.
}

%%%%%%%%%%%%%%%%%%%%%%%%%%%%%%%%%%%%%%%%%%%%%%%%%%%%%%%%%%%%
% SECTION SOURCE TRACKING
%%%%%%%%%%%%%%%%%%%%%%%%%%%%%%%%%%%%%%%%%%%%%%%%%%%%%%%%%%%%

%------------------------------------------------------------
% 14.14 Statistical Tables
%------------------------------------------------------------
%
% Source basis:
%   - Standard descriptive and inferential statistics.
%   - Standard regression, ANOVA, categorical-data, and
%     statistical testing formulas.
%
% Used for:
%   - population and sample notation;
%   - means;
%   - weighted means;
%   - variances and standard deviations;
%   - ranges and quartiles;
%   - z-scores;
%   - proportions;
%   - covariance;
%   - correlation;
%   - standard errors;
%   - finite-population correction;
%   - critical values;
%   - confidence intervals;
%   - z and t tests;
%   - proportion tests;
%   - p-values;
%   - Type I and Type II errors;
%   - power;
%   - chi-square tests;
%   - contingency tables;
%   - regression;
%   - residuals;
%   - sums of squares;
%   - R-squared;
%   - ANOVA;
%   - odds, odds ratios, relative risks;
%   - effect sizes;
%   - multiple-testing corrections;
%   - common errors;
%   - applications and exercises.
%
% Notes:
%   - Common Mathematical Theorems is developed next in Section 14.15.
%   - Mathematical Terminology follows in Section 14.16.
%
%%%%%%%%%%%%%%%%%%%%%%%%%%%%%%%%%%%%%%%%%%%%%%%%%%%%%%%%%%%%